% META: {"id": "TSPE-SUITLIM-CO-015", "chapitre": "TSPE-SUITES-LIMITES", "type_objet": "corrige", "exercice_ref": "TSPE-SUITLIM-EX-015", "capacites_codes": ["C3"], "sources_inspiration": [], "status": "generated"}
% BEGIN-VERIFY
% from sympy import *
% n = symbols('n', integer=True, nonnegative=True)
% a = Rational(4, 5)
% u_expr = 250 - 150 * a**n
% assert u_expr.subs(n, 0) == 100
% assert u_expr.subs(n, 1) == 130
% assert limit(u_expr, n, oo) == 250
% END-VERIFY
\begin{corrige}{TSPE-SUITLIM-EX-015}

\textbf{Question 1.} $u_0 = 100$ (volume initial). Chaque jour, il reste $80\,\%$ de l'eau plus $50$~L : $u_{n+1} = 0{,}8\,u_n + 50$.

\textbf{Question 2.} $u_1 = 0{,}8 \times 100 + 50 = 130$. $u_2 = 0{,}8 \times 130 + 50 = 154$.

\textbf{Question 3.} $\ell = 0{,}8\ell + 50 \Rightarrow 0{,}2\ell = 50 \Rightarrow \ell = 250$.

$v_n = u_n - 250$, $v_{n+1} = 0{,}8\,u_n + 50 - 250 = 0{,}8(u_n - 250) = 0{,}8\,v_n$. $(v_n)$ geometrique de raison $0{,}8$.

\textbf{Question 4.} $v_n = (100-250) \times 0{,}8^n = -150 \times 0{,}8^n$. $u_n = 250 - 150 \times 0{,}8^n$.

$0{,}8^n \to 0$, donc $\boxed{\lim u_n = 250}$~litres.

\end{corrige}
